\documentclass[a4paper,12pt]{article}
\usepackage{amsmath,amssymb}
\usepackage{geometry}
\geometry{left=2cm, right=2cm, top=2.5cm, bottom=2.5cm}
\usepackage{xcolor}
\usepackage[most]{tcolorbox}
\usepackage{cancel}

% Define colors
\definecolor{SolutionColor}{HTML}{F0F8FF}
\definecolor{ExerciseColor}{HTML}{E6E6FA}

% Define environments
\newtcolorbox{solutionbox}[1]{
    colback=SolutionColor,
    colframe=blue!75!black,
    title=#1,
    fonttitle=\bfseries,
    breakable
}

\newtcolorbox{exercice}[1][]{
    colback=ExerciseColor,
    colframe=purple!75!black,
    title=#1,
    fonttitle=\bfseries,
    breakable
}

\begin{document}

% Exercice 9
\begin{exercice}[title=Exercice 9]
Factoriser les expressions suivantes puis résoudre l’équation \( f(x) = 0 \).

\begin{enumerate}
    \item \( f(x) = x^2 + x \)
    \item \( f(x) = 3x^2 + 2x \)
    \item \( f(x) = 5x^2 - 7x \)
    \item \( f(x) = x(2x - 1) + x(5x - 2) \)
    \item \( f(x) = 2x(2x + 5) - 3(2x - 2)x \)
    \item \( f(x) = (x + 3)(2x + 1) + (x + 3)(x - 4) \)
    \item \( f(x) = (2x - 5)(2x + 6) - (2x - 5)(x + 4) \)
    \item \( f(x) = (3x + 1)(x - 4) - (3x + 1) \)
    \item \( f(x) = 2(2x - 1)(3x + 5) + (3x + 5)(x - 1) \)
    \item \( f(x) = (x + 3)(2x - 5) + (2x + 6)(x - 2) \)
    \item \( f(x) = (2x - 4)^2 - (2x - 4)(x + 5) \)
    \item \( f(x) = (x - 2)(x + 1) - (x - 2)(x - 4) + 3(x + 2)(x - 2) \)
    \item \( f(x) = (x + 3)^2 - 25 \)
    \item \( f(x) = (3x - 4)^2 - (x + 1)^2 \)
    \item \( f(x) = x^2 - 6x + 9 \)
    \item \( f(x) = 4x^2 + 20x + 25 \)
\end{enumerate}
\end{exercice}

\begin{solutionbox}{Solution}
\begin{enumerate}
    \item \( f(x) = x^2 + x \)

    Factoriser \( f(x) \) :
    \[
    f(x) = x(x + 1)
    \]

    Résoudre \( f(x) = 0 \) :
    \[
    x(x + 1) = 0 \quad \Rightarrow \quad x = 0 \quad \text{ou} \quad x + 1 = 0 \quad \Rightarrow \quad x = -1
    \]

    \textbf{Solution :} \( x = 0 \) ou \( x = -1 \).

    \item \( f(x) = 3x^2 + 2x \)

    Factoriser :
    \[
    f(x) = x(3x + 2)
    \]

    Résoudre \( f(x) = 0 \) :
    \[
    x(3x + 2) = 0 \quad \Rightarrow \quad x = 0 \quad \text{ou} \quad 3x + 2 = 0 \quad \Rightarrow \quad x = -\dfrac{2}{3}
    \]

    \textbf{Solution :} \( x = 0 \) ou \( x = -\dfrac{2}{3} \).

    \item \( f(x) = 5x^2 - 7x \)

    Factoriser :
    \[
    f(x) = x(5x - 7)
    \]

    Résoudre \( f(x) = 0 \) :
    \[
    x(5x - 7) = 0 \quad \Rightarrow \quad x = 0 \quad \text{ou} \quad 5x - 7 = 0 \quad \Rightarrow \quad x = \dfrac{7}{5}
    \]

    \textbf{Solution :} \( x = 0 \) ou \( x = \dfrac{7}{5} \).

    \item \( f(x) = x(2x - 1) + x(5x - 2) \)

    Factoriser :
    \[
    f(x) = x[(2x - 1) + (5x - 2)] = x(7x - 3)
    \]

    Résoudre \( f(x) = 0 \) :
    \[
    x(7x - 3) = 0 \quad \Rightarrow \quad x = 0 \quad \text{ou} \quad 7x - 3 = 0 \quad \Rightarrow \quad x = \dfrac{3}{7}
    \]

    \textbf{Solution :} \( x = 0 \) ou \( x = \dfrac{3}{7} \).

    \item \( f(x) = 2x(2x + 5) - 3(2x - 2)x \)

    Développons et simplifions :
    \[
    f(x) = [4x^2 + 10x] - [6x^2 - 6x] = -2x^2 + 16x
    \]

    Factorisons :
    \[
    f(x) = -2x(x - 8)
    \]

    Résoudre \( f(x) = 0 \) :
    \[
    -2x(x - 8) = 0 \quad \Rightarrow \quad x = 0 \quad \text{ou} \quad x - 8 = 0 \quad \Rightarrow \quad x = 8
    \]

    \textbf{Solution :} \( x = 0 \) ou \( x = 8 \).

    \item \( f(x) = (x + 3)(2x + 1) + (x + 3)(x - 4) \)

    Factorisons \( (x + 3) \) :
    \[
    f(x) = (x + 3)[(2x + 1) + (x - 4)] = (x + 3)(3x - 3) = 3(x + 3)(x - 1)
    \]

    Résoudre \( f(x) = 0 \) :
    \[
    3(x + 3)(x - 1) = 0 \quad \Rightarrow \quad x = -3 \quad \text{ou} \quad x = 1
    \]

    \textbf{Solution :} \( x = -3 \) ou \( x = 1 \).

    \item \( f(x) = (2x - 5)(2x + 6) - (2x - 5)(x + 4) \)

    Factorisons \( (2x - 5) \) :
    \[
    f(x) = (2x - 5)[(2x + 6) - (x + 4)] = (2x - 5)(x + 2)
    \]

    Résoudre \( f(x) = 0 \) :
    \[
    (2x - 5)(x + 2) = 0 \quad \Rightarrow \quad x = \dfrac{5}{2} \quad \text{ou} \quad x = -2
    \]

    \textbf{Solution :} \( x = \dfrac{5}{2} \) ou \( x = -2 \).

    \item \( f(x) = (3x + 1)(x - 4) - (3x + 1) \)

    Factorisons \( (3x + 1) \) :
    \[
    f(x) = (3x + 1)[(x - 4) - 1] = (3x + 1)(x - 5)
    \]

    Résoudre \( f(x) = 0 \) :
    \[
    (3x + 1)(x - 5) = 0 \quad \Rightarrow \quad x = -\dfrac{1}{3} \quad \text{ou} \quad x = 5
    \]

    \textbf{Solution :} \( x = -\dfrac{1}{3} \) ou \( x = 5 \).

    \item \( f(x) = 2(2x - 1)(3x + 5) + (3x + 5)(x - 1) \)

    Factorisons \( (3x + 5) \) :
    \[
    f(x) = (3x + 5)[2(2x - 1) + (x - 1)] = (3x + 5)(5x - 3)
    \]

    Résoudre \( f(x) = 0 \) :
    \[
    (3x + 5)(5x - 3) = 0 \quad \Rightarrow \quad x = -\dfrac{5}{3} \quad \text{ou} \quad x = \dfrac{3}{5}
    \]

    \textbf{Solution :} \( x = -\dfrac{5}{3} \) ou \( x = \dfrac{3}{5} \).

    \item \( f(x) = (x + 3)(2x - 5) + (2x + 6)(x - 2) \)

    Remarquons que \( 2x + 6 = 2(x + 3) \).

    Factorisons \( (x + 3) \) :
    \[
    f(x) = (x + 3)(2x - 5) + 2(x + 3)(x - 2) = (x + 3)[(2x - 5) + 2(x - 2)] = (x + 3)(4x - 9)
    \]

    Résoudre \( f(x) = 0 \) :
    \[
    (x + 3)(4x - 9) = 0 \quad \Rightarrow \quad x = -3 \quad \text{ou} \quad x = \dfrac{9}{4}
    \]

    \textbf{Solution :} \( x = -3 \) ou \( x = \dfrac{9}{4} \).

    \item \( f(x) = (2x - 4)^2 - (2x - 4)(x + 5) \)

    Factorisons \( (2x - 4) \) :
    \[
    f(x) = (2x - 4)[(2x - 4) - (x + 5)] = (2x - 4)(x - 9)
    \]

    Factorisons \( 2x - 4 = 2(x - 2) \) :
    \[
    f(x) = 2(x - 2)(x - 9)
    \]

    Résoudre \( f(x) = 0 \) :
    \[
    2(x - 2)(x - 9) = 0 \quad \Rightarrow \quad x = 2 \quad \text{ou} \quad x = 9
    \]

    \textbf{Solution :} \( x = 2 \) ou \( x = 9 \).

    \item \( f(x) = (x - 2)(x + 1) - (x - 2)(x - 4) + 3(x + 2)(x - 2) \)

    Factorisons \( (x - 2) \) :
    \[
    f(x) = (x - 2)\left[ (x + 1) - (x - 4) + 3(x + 2) \right]
    \]
    Simplifions :
    \[
    (x + 1) - (x - 4) = x + 1 - x + 4 = 5
    \]
    \[
    f(x) = (x - 2)[5 + 3(x + 2)] = (x - 2)[5 + 3x + 6] = (x - 2)(3x + 11)
    \]

    Résoudre \( f(x) = 0 \) :
    \[
    (x - 2)(3x + 11) = 0 \quad \Rightarrow \quad x = 2 \quad \text{ou} \quad x = -\dfrac{11}{3}
    \]

    \textbf{Solution :} \( x = 2 \) ou \( x = -\dfrac{11}{3} \).

    \item \( f(x) = (x + 3)^2 - 25 \)

    Reconnaissons une différence de carrés :
    \[
    f(x) = (x + 3 - 5)(x + 3 + 5) = (x - 2)(x + 8)
    \]

    Résoudre \( f(x) = 0 \) :
    \[
    (x - 2)(x + 8) = 0 \quad \Rightarrow \quad x = 2 \quad \text{ou} \quad x = -8
    \]

    \textbf{Solution :} \( x = 2 \) ou \( x = -8 \).

    \item \( f(x) = (3x - 4)^2 - (x + 1)^2 \)

    Utilisons l'identité remarquable \( a^2 - b^2 = (a - b)(a + b) \) :
    \[
    f(x) = [(3x - 4) - (x + 1)][(3x - 4) + (x + 1)] = (2x - 5)(4x - 3)
    \]

    Résoudre \( f(x) = 0 \) :
    \[
    (2x - 5)(4x - 3) = 0 \quad \Rightarrow \quad x = \dfrac{5}{2} \quad \text{ou} \quad x = \dfrac{3}{4}
    \]

    \textbf{Solution :} \( x = \dfrac{5}{2} \) ou \( x = \dfrac{3}{4} \).

    \item \( f(x) = x^2 - 6x + 9 \)

    Reconnaissons un carré parfait :
    \[
    f(x) = (x - 3)^2
    \]

    Résoudre \( f(x) = 0 \) :
    \[
    (x - 3)^2 = 0 \quad \Rightarrow \quad x = 3
    \]

    \textbf{Solution :} \( x = 3 \).

    \item \( f(x) = 4x^2 + 20x + 25 \)

    Reconnaissons un carré parfait :
    \[
    f(x) = (2x + 5)^2
    \]

    Résoudre \( f(x) = 0 \) :
    \[
    (2x + 5)^2 = 0 \quad \Rightarrow \quad x = -\dfrac{5}{2}
    \]

    \textbf{Solution :} \( x = -\dfrac{5}{2} \).

\end{enumerate}
\end{solutionbox}

% Exercice 10
\begin{exercice}[title=Exercice 10]
Résoudre les équations suivantes :

\begin{enumerate}
    \item \( x^2 = 9 \)
    \item \( x^2 + 1 = 17 \)
    \item \( 2x^2 - 4 = 2 \)
    \item \( 5x^2 - 4 = -4 \)
    \item \( -4x^2 - 2 = 10 \)
    \item \( (2x + 1)(x^2 - 16) = 0 \)
    \item \( (3x + 6)(x^2 - 2) = 0 \)
    \item \( x(x^2 + 3) = 0 \)
    \item \( 2x^3 + 5x = 0 \)
    \item \( x^4 - x^2 = 0 \)
\end{enumerate}
\end{exercice}

\begin{solutionbox}{Solution}
\begin{enumerate}
    \item \( x^2 = 9 \)

    \[
    x = \pm 3
    \]

    \textbf{Solution :} \( x = -3 \) ou \( x = 3 \).

    \item \( x^2 + 1 = 17 \)

    \[
    x^2 = 16 \quad \Rightarrow \quad x = \pm 4
    \]

    \textbf{Solution :} \( x = -4 \) ou \( x = 4 \).

    \item \( 2x^2 - 4 = 2 \)

    \[
    2x^2 = 6 \quad \Rightarrow \quad x^2 = 3 \quad \Rightarrow \quad x = \pm \sqrt{3}
    \]

    \textbf{Solution :} \( x = -\sqrt{3} \) ou \( x = \sqrt{3} \).

    \item \( 5x^2 - 4 = -4 \)

    \[
    5x^2 = 0 \quad \Rightarrow \quad x = 0
    \]

    \textbf{Solution :} \( x = 0 \).

    \item \( -4x^2 - 2 = 10 \)

    \[
    -4x^2 = 12 \quad \Rightarrow \quad x^2 = -3
    \]

    Pas de solution réelle.

    \textbf{Solution :} Aucune solution réelle.

    \item \( (2x + 1)(x^2 - 16) = 0 \)

    \[
    2x + 1 = 0 \quad \Rightarrow \quad x = -\dfrac{1}{2} \\
    x^2 - 16 = 0 \quad \Rightarrow \quad x = \pm 4
    \]

    \textbf{Solution :} \( x = -\dfrac{1}{2},\ x = -4,\ x = 4 \).

    \item \( (3x + 6)(x^2 - 2) = 0 \)

    \[
    3x + 6 = 0 \quad \Rightarrow \quad x = -2 \\
    x^2 - 2 = 0 \quad \Rightarrow \quad x = \pm \sqrt{2}
    \]

    \textbf{Solution :} \( x = -2,\ x = -\sqrt{2},\ x = \sqrt{2} \).

    \item \( x(x^2 + 3) = 0 \)

    \[
    x = 0 \quad \text{ou} \quad x^2 + 3 = 0 \quad (\text{Pas de solution réelle})
    \]

    \textbf{Solution :} \( x = 0 \).

    \item \( 2x^3 + 5x = 0 \)

    Factorisons :
    \[
    x(2x^2 + 5) = 0
    \]

    \[
    x = 0 \quad \text{ou} \quad 2x^2 + 5 = 0 \quad (\text{Pas de solution réelle})
    \]

    \textbf{Solution :} \( x = 0 \).

    \item \( x^4 - x^2 = 0 \)

    Factorisons :
    \[
    x^2(x^2 - 1) = 0 \quad \Rightarrow \quad x^2 = 0 \quad \text{ou} \quad x^2 - 1 = 0
    \]

    \[
    x = 0 \quad \text{ou} \quad x = \pm 1
    \]

    \textbf{Solution :} \( x = -1,\ x = 0,\ x = 1 \).

\end{enumerate}
\end{solutionbox}

% Exercice 11
\begin{exercice}[title=Exercice 11]
Simplifier les fractions suivantes :

\begin{enumerate}
    \item \( \dfrac{5(x+3)}{2(x+3)} \)
    \item \( \dfrac{3x+9}{x+1} \)
    \item \( \dfrac{x^2 - 1}{x + 1} \)
    \item \( \dfrac{x^2 - 6x + 9}{x^2 - 9} \)
    \item \( \dfrac{x^2 + 4x + 4}{5x + 10} \)
    \item \( \dfrac{3(x+5) + 2(x+5)}{4(x+5)} \)
    \item \( \dfrac{x(x - 4) + 2x - 8}{x - 4} \)
    \item \( \dfrac{x(x^2 - 4)}{x^2 - 2x} \)
    \item \( \dfrac{3(x - 1)(x^2 - 2)}{(x - \sqrt{2})(x^2 + 1)} \)
\end{enumerate}
\end{exercice}

\begin{solutionbox}{Solution}
\begin{enumerate}
    \item \( \dfrac{5(x+3)}{2(x+3)} \)

    Simplifions :
    \[
    \dfrac{5\cancel{(x + 3)}}{2\cancel{(x + 3)}} = \dfrac{5}{2}
    \]

    \textbf{Forme simplifiée :} \( \dfrac{5}{2} \).

    \item \( \dfrac{3x+9}{x+1} \)

    Factorisons le numérateur :
    \[
    \dfrac{3(x + 3)}{x + 1}
    \]

    Pas de simplification possible.

    \textbf{Forme simplifiée :} \( \dfrac{3(x + 3)}{x + 1} \).

    \item \( \dfrac{x^2 - 1}{x + 1} \)

    Reconnaissons une différence de carrés :
    \[
    \dfrac{(x - 1)(x + 1)}{x + 1} = \dfrac{\cancel{(x + 1)}(x - 1)}{\cancel{(x + 1)}} = x - 1
    \]

    \textbf{Forme simplifiée :} \( x - 1 \).

    \item \( \dfrac{x^2 - 6x + 9}{x^2 - 9} \)

    Factorisons :
    \[
    \dfrac{(x - 3)^2}{(x - 3)(x + 3)} = \dfrac{\cancel{(x - 3)}(x - 3)}{\cancel{(x - 3)}(x + 3)} = \dfrac{x - 3}{x + 3}
    \]

    \textbf{Forme simplifiée :} \( \dfrac{x - 3}{x + 3} \).

    \item \( \dfrac{x^2 + 4x + 4}{5x + 10} \)

    Factorisons :
    \[
    \dfrac{(x + 2)^2}{5(x + 2)} = \dfrac{\cancel{(x + 2)}(x + 2)}{5\cancel{(x + 2)}} = \dfrac{x + 2}{5}
    \]

    \textbf{Forme simplifiée :} \( \dfrac{x + 2}{5} \).

    \item \( \dfrac{3(x+5) + 2(x+5)}{4(x+5)} \)

    Simplifions le numérateur :
    \[
    \dfrac{5(x + 5)}{4(x + 5)} = \dfrac{5}{4}
    \]

    \textbf{Forme simplifiée :} \( \dfrac{5}{4} \).

    \item \( \dfrac{x(x - 4) + 2x - 8}{x - 4} \)

    Développons et simplifions :
    \[
    \dfrac{x^2 - 4x + 2x - 8}{x - 4} = \dfrac{x^2 - 2x - 8}{x - 4}
    \]
    Factorisons le numérateur :
    \[
    \dfrac{(x - 4)(x + 2)}{x - 4} = x + 2
    \]

    \textbf{Forme simplifiée :} \( x + 2 \).

    \item \( \dfrac{x(x^2 - 4)}{x^2 - 2x} \)

    Factorisons :
    \[
    \dfrac{x(x - 2)(x + 2)}{x(x - 2)} = x + 2
    \]

    \textbf{Forme simplifiée :} \( x + 2 \).

    \item \( \dfrac{3(x - 1)(x^2 - 2)}{(x - \sqrt{2})(x^2 + 1)} \)

    Factorisons \( x^2 - 2 = (x - \sqrt{2})(x + \sqrt{2}) \) :
    \[
    \dfrac{3(x - 1)(x - \sqrt{2})(x + \sqrt{2})}{(x - \sqrt{2})(x^2 + 1)} = \dfrac{3(x - 1)(x + \sqrt{2})}{x^2 + 1}
    \]

    \textbf{Forme simplifiée :} \( \dfrac{3(x - 1)(x + \sqrt{2})}{x^2 + 1} \).

\end{enumerate}
\end{solutionbox}

\end{document}